% 《大学物理学学习指导》自测题11计算11.3.1—11.3.6
% 原书印刷232，PDF物理240页。源题数据、单位与边界逐项核对。
\subsection{自测题11：光电、康普顿、氢谱和势阱计算题}

\sourcetest{11.3.1}{232}{240}
\problemtext 410nm单色光照射金属，最大光电子动能为1.0eV，
求能使该金属发生光电效应的单色光最大波长。
\workings
由爱因斯坦光电方程
$W=hc/\lambda-K_{\max}$，采用
$hc\simeq1239.842\,\mathrm{eV\,nm}$，
\[
 W=\frac{1239.842}{410}-1.0
 \simeq2.024\,\mathrm{eV}.
\]
红限波长由$hc/\lambda_0=W$给出，
\[
\boxed{\lambda_{\max}=\lambda_0
=\frac{1239.842}{1239.842/410-1}\,\mathrm{nm}
\simeq612.6\,\mathrm{nm}}.
\]
\crosscheck 红限波长应\emph{长于}410nm，
因为原入射光子有1eV的多余动能可供电子带走。

\sourcetest{11.3.2}{232}{240}
\problemtext 0.0500nm的X射线在自由静止电子上发生康普顿散射；
散射后波长0.0522nm。
求光子的散射角及电子得到的动能。
\workings
康普顿波长
$\lambda_C=h/(m_ec)=0.00242631$nm。
波长增加
$\Delta\lambda=0.0522-0.0500=0.0022$nm，
\[
\cos\theta=1-\frac{\Delta\lambda}{\lambda_C}
=1-\frac{0.0022}{0.00242631}
\simeq0.09328,
\]
因此
\[
\boxed{\theta\simeq84.65^\circ}.
\]
电子初始静止，反冲动能为光子损失的能量：
\[
\begin{aligned}
K_e&=hc\left(\frac1{\lambda_0}-\frac1{\lambda'}\right)\\
&\simeq(1239.842\,\mathrm{eV\,nm})
\left(\frac1{0.0500}-\frac1{0.0522}\right).
\end{aligned}
\]
从而
\[
\boxed{K_e\simeq1.045\times10^3\,\mathrm{eV}
=1.045\,\mathrm{keV}}.
\]
\crosscheck $\lambda'>\lambda_0$，故电子得到正动能；
将$\lambda,\lambda'$倒置会产生负能量，违反题设。

\sourcetest{11.3.3}{232}{240}
\problemtext 氢原子电子从基态跃迁至主量子数$n$的激发态，
这时玻尔轨道半径变为原来9倍，求原子从该态
直接返回基态时发射光子的频率。
\workings
玻尔轨道半径$r_n=n^2a_0$，
$r_n/r_1=9$给$n=3$。
氢能级$E_n=-13.6/n^2$eV，
$n=3\to1$光子能量
\[
\Delta E=13.6\left(1-\frac1{9}\right)
\simeq12.0889\,\mathrm{eV}.
\]
因此
\[
\boxed{\nu=\frac{\Delta E}{h}
=\frac{12.0889\,e}{h}
\simeq2.923\times10^{15}\,\mathrm{Hz}}.
\]
\crosscheck 题中“基态跃迁到激发态”本身是\emph{吸收}能量；
要谈\emph{发射}光子需明确从$n=3$返回较低能级，
此处按源题的直接返回基态模型给出频率。

\sourcetest{11.3.4}{232}{240}
\problemtext 氢原子从某个初始态跃迁到
\emph{从基态激发需10.19eV的能级}时
发射波长486nm的光子，
求初始态的能量与主量子数。
\workings
从基态$E_1=-13.6$eV激发10.19eV，
终态能量约$E_f=-3.41$eV，
即按玻尔能级取$n_f=2$、$E_f\simeq-3.4$eV。
光子能量
$E_\gamma=hc/(486\,\mathrm{nm})\simeq2.551$eV。
于是初始电子能量
\[
E_i=E_f+E_\gamma\simeq-3.4+2.55
=-0.85\,\mathrm{eV}.
\]
解$E_i=-13.6/n_i^2$，
\[
\boxed{n_i=4},\qquad
\boxed{E_i=-0.85\,\mathrm{eV}}.
\]
\crosscheck 巴耳末线486nm对应$4\to2$；
不能因为题干出现“从基态激发能10.19eV”
就错把486nm理解成$1\to2$发射。

\sourcetest{11.3.5}{232}{240}
\problemtext 若采用能量等效质量$m_{\rm eq}=E/c^2$的旧教材表述，
求波长550nm光子的等效质量（kg）。
\workings
光子能量$E=hc/\lambda$，所以
\[
\boxed{m_{\rm eq}=\frac{E}{c^2}
=\frac{h}{\lambda c}
=\frac{6.62607015\times10^{-34}}
{(550\times10^{-9})(299792458)}
\simeq4.02\times10^{-36}\,\mathrm{kg}}.
\]
\crosscheck 这不是光子的\emph{静质量}；
现代相对论表述中光子静质量为零，
这里的$m_{\rm eq}$仅是光子能量除以$c^2$。

\sourcetest{11.3.6}{232}{240}
\problemtext 一维无限深势阱$0\le x\le a$内，
已归一化波函数为
$\psi(x)=\sqrt{2/a}\sin(\pi x/a)$，
求找到粒子概率密度最大的位置。
\workings
概率密度
\[
\rho(x)=|\psi(x)|^2
=\frac2a\sin^2\left(\frac{\pi x}{a}\right).
\]
在$[0,a]$范围内，
$\sin^2(\pi x/a)$于$\pi x/a=\pi/2$处达到最大1，
故
\[
\boxed{x=\frac a2},\qquad
\boxed{\rho_{\max}=\frac2a}.
\]
\crosscheck 势阱两端$x=0,a$的波函数为零，
故概率密度最大的位置在阱的正中心，而不是两端。
